Let $C^+$ be the fiber of ordinary global cochains to ordinary real
cochains. Real $H^0$ is included. Over a finite Artin quotient,
continuous cochains are filtered unions of finite free function
modules, hence flat and therefore free. A residual splitting into
finite cohomology and disks lifts by the nilpotent perturbation
series to a finite free model with actual homotopies. Minimal models
at successive quotients are related by invertible matrices:
a homotopy equivalence reduces to an isomorphism of minimal residual
complexes, and Nakayama applies. Lifting these matrices and comparison
homotopies makes the models compatible. Their acyclic, termwise
surjective comparison cones have acyclic inverse limit.
Here is the compatibility argument. At successive finite coefficient
quotients, reduction of a minimal finite free model is a minimal model
of the reduced cochains. Two such models are homotopy equivalent;
reducing a homotopy inverse modulo the maximal ideal shows that the
comparison is degreewise invertible. Lift these invertible matrices
to the next quotient and change bases. The model differentials then
reduce to the preceding differentials exactly. The comparison maps
to the original cochains agree after reduction up to chain homotopy.
Since reduction of continuous cochains is termwise surjective, lift
that homotopy and adjust the upper comparison by its boundary. The
comparison maps are now exactly compatible and remain
quasi-isomorphisms.

Let $\mathcal K_n$ be their acyclic comparison cones. Their transition
maps are termwise surjective. If $s$ denotes the transition shift on
the product, successive lifting gives the exact sequence of complexes
\[
 0\longrightarrow\varprojlim_n\mathcal K_n
 \longrightarrow\prod_n\mathcal K_n
 \xrightarrow{\,1-s\,}\prod_n\mathcal K_n
 \longrightarrow0.
\]
Products of complexes of modules preserve acyclicity, so the long
exact cohomology sequence makes $\varprojlim_n\mathcal K_n$ acyclic.
The inverse-limit comparison is therefore a quasi-isomorphism from
the bounded finite free limit model to the continuous coefficient
complex. This also identifies its derived coefficient changes with
the finite-quotient constructions.
This proves the needed perfection and coefficient change.

% Replace baseline lines 3310--3335; retains equation tag (E.3).
